% Appendix: the problem statement in detail. Moved from Section~\ref{sec:problem}; numbers unchanged.

\section{The problem statement in detail}\label{app:problem}

This appendix holds the discussion that Section~\ref{sec:problem} summarizes: how the criteria relate to their
sources, what follows from the observation, the cache and count arguments with the measured delays, and the full
threat model and metrics.

\subsection{Reading the criteria}\label{app:criteria}

Security, and freshness for the $\delta$ that the serving path states, are required of every result,
and count-honesty is required of the flag. Soundness and completeness are what approximate search
trades away; permission-recall@$k$, the under-retrieval rate and the false-completeness rate
(Section~\ref{sec:prob-metrics}) measure the shortfall. Count-honesty is relative to $\varphi_s$: it
implies completeness only when $\varphi_s$ has no false denials, $R_\tau(s) \subseteq E(\varphi_s)$,
and it never implies soundness, since a full answer of the wrong chunks satisfies it. A flag with
$c = 1$ only for maximal results would need exact search to set; Appendix~\ref{sec:sys-complete} shows
that the layer we test is count-honest relative to $\varphi_s$ and no more. Wang et al.'s soundness
compares an answer with the unrestricted one; ours compares it with the exact answer over the
readable set, which is what a Truman subject is owed. Score fidelity is a separate conformance
property (finite scores in non-increasing order, Section~\ref{sec:conformance}); mis-ordering inside
a store, such as the slightly out-of-order output of pgvector's relaxed iterative
scans~\cite{pgvector2026pgvector}, shows in recall, not in security. Security concerns the returned
set only. B\"{u}ttcher and Clarke's stricter requirement,
that no aspect of the output depend on unsearchable content~\cite{buttcher2005security}, is a
noninterference condition~\cite{goguen1982security} on everything a caller observes, including timing
and the flag; we treat those observables as side channels (Section~\ref{sec:prob-threat}). Under
Truman semantics a short answer looks like an answer to a subject who can read little; the flag $c$
is a partial step toward non-Truman behavior, letting the layer say that an answer may be
incomplete.

With the readable set defined by the policy and (A1) equating the check with the policy, the
statement is close to definitional, and we call it an observation for that reason. Its content is
what it leaves out: no assumption about $\varphi_s$, the filter APIs, the translators, the store's
filter semantics or its index. It moves the burden of security onto a small trusted computing base,
stated in Section~\ref{sec:prob-criteria}, and turns everything else into a question of completeness, soundness and freshness. It
also holds for every query, chosen adaptively or not, which no sampled measurement can show.

Three consequences follow. \emph{First, the re-check sees filter faults in one direction only.} A
\emph{false admit}, where $\varphi_s$ accepts a chunk outside $R_\tau(s)$, reaches the re-check, which
drops it and counts a violation; the violation count thus measures online how far the pushed-down
filter over-approximates. A \emph{false denial}, where $\varphi_s$ rejects a readable chunk, never
becomes a candidate and silently costs completeness. Truncated allow-lists (OpenFGA's non-streaming
ListObjects stops at 1,000 results or 3 seconds by default~\cite{openfga2026openfga}), divergent store
filter semantics and approximate-index misses all fail this way, and only an offline oracle finds
them: differential conformance testing against a reference semantics
(Section~\ref{sec:conformance}) or an audit against exact search over $R_\tau(s)$
(Section~\ref{sec:quality}). \emph{Second, the assumptions mark where security can still fail.} (A1)
excludes wrong upstream permissions and, when the check reads ACLs copied into chunk metadata, holds
only as far as the copy is fresh; (A2) fails under any store or ingestion fault that binds content to
the wrong identifier or metadata, which no re-check can detect; and security holds at the revision
the check reads, so how stale that revision may be is a separate guarantee, freshness, which we
state per path below. \emph{Third, a violation is a translation fault only at a single revision.} If
$\varphi_s$ is compiled at revision $\tau_f$ and hits are checked at $\tau_c \neq \tau_f$, a
disagreement may be a legitimate ACL change between the calls. Zanzibar's snapshot tokens and
SpiceDB's revision tokens prevent such mixing~\cite{pang2019zanzibar,authzed2026protecting}, and we pin
filter and check to one revision when measuring disagreement (Section~\ref{sec:engines}).

\paragraph{The trusted computing base in full.}
The observation trusts (i)~the authorizer's check API and the authorization data it reads (A1),
including the freshness of ACLs copied into chunk metadata; (ii)~the layer's mapping of a request to
checks: the subject taken from a verified token (Appendix~\ref{sec:sys-identity}) and the adapter code
that turns the subject, its groups and each chunk's identifier and attributes into the engine's check
request; (iii)~the binding of content to identifier and metadata in the store and in ingestion (A2);
(iv)~the decision cache, its keys, time to live and invalidation, and the engine's own staleness,
which together set $\delta$; and (v)~the fail-closed path (A3). It does not trust the filter APIs, the
translators, the store's filter semantics or index, the widening loop, or a reranker or redactor,
which receive only permitted hits. The trusted part has failed before. Two faults in the layer's own
adapters let the check itself admit, beyond any re-check: a SpiceDB adapter that checked groups as
separate subjects and so bypassed an exclusion, and a Cedar adapter that dropped values a
\texttt{forbid} rule needed (Table~\ref{tab:conf-faults}); the project's acceptance scenarios and
review found them, and both were fixed before the version we evaluate. With the re-check turned off,
the filter and its translator join the trusted computing base (Appendix~\ref{sec:sys-retrieval}).

\subsection{Caches, staleness and count certificates}\label{app:certificates}

\paragraph{Caches and staleness.}
Let the engine answer from a snapshot that may lack the writes of the last $\delta_e$ seconds, and let
the layer cache each decision for at most a time to live $\ell$. A decision used at time $t$ was then
read at a revision that includes every write committed before $t - (\ell + \delta_e)$, so under (A1)
to (A3) the layer is $(\ell + \delta_e)$-fresh, or $\delta_e$-fresh without a cache. An explicit
invalidation is meant to restore $\delta_e$-freshness for the subjects it names; it can do so only if
the cache does not then store answers the engine computes from a snapshot older than the
invalidation. Section~\ref{sec:eng-revoke} measures $\delta$ per path (Table~\ref{tab:eng-revoke}).
When permissions are copied into chunk metadata, and for engines read at full consistency, revocation
took effect in a median of 0.01\,s or less; for engines behind a cache it took the time to live
(medians within 34\,ms of 5\,s and 30\,s), and invalidation ended that delay at once, or within 1\,s
on a replica sharing the cache. Two sources of truth were stale themselves. Under SpiceDB's default
\texttt{minimizeLatency}, revoked chunks were served for a median 3.3\,s and, in one of 40 trials of
the first run, 8.7\,s, beyond the 6.5\,s we had registered as the limit (the documented 5.5\,s plus
1\,s for polling); a rerun stayed within 5.2\,s. Cerbos took a median 2.5\,s to reload a rewritten policy file, and
states no bound. In front of these two engines, the release we evaluated violated its stated bound:
after an invalidation the next call refilled the cache from the stale snapshot and held it for the
full time to live, so 79 of 80 SpiceDB revocations with invalidation exceeded the bound, with medians
of 30 and 31\,s. The fix caches nothing for the engine's declared staleness after an invalidation, and
a rerun confirmed it: all 160 SpiceDB and Cerbos revocations with invalidation took effect within the
declared staleness plus 1\,s, and 1\,s more on the replica.

\paragraph{Count certificates.}
Security needs no knowledge of $R_\tau(s)$; completeness does. A layer can set $c = 1$ count-honestly
only if (a)~the store searched exactly under $\varphi_s$, deep enough that $k$ candidates survived the
re-check or none remained, or (b)~a filter-only count shows that every chunk passing $\varphi_s$ was
among the candidates. Conversely, a count of at least $k$ alongside fewer than $k$ candidates proves
under-retrieval; OpenSearch uses this test to switch to exact search~\cite{opensearch2026efficient}.
Both certificates are relative to $\varphi_s$, so completeness with respect to $R_\tau(s)$ also needs
a filter free of false denials, which we test offline rather than certify at run time. A layer that
marks a short or empty approximate result complete is not count-honest; the release of \sys{} we ran
end to end was not, and in the runs of Section~\ref{sec:e2e} it reported empty answers as complete for
up to ten least-privileged people per run, each of whom could read 624 chunks.

\subsection{The threat model in full}\label{app:threat}

\paragraph{Adversary.}
The adversary is an authenticated subject, or an agent or application acting for one, who may issue
any number of arbitrary queries through the layer, chosen adaptively and including adversarial and
red-team prompts, to obtain the content of any chunk outside the subject's readable set, across
users or tenants, or to learn whether such chunks exist near a query. The adversary does not control
the authorizer, the store or the layer's host. Observation~\ref{prop:reduction} covers adaptive
queries; our measurements sample questions from a generator and bound failure rates over that
distribution only (Section~\ref{sec:prob-metrics}).

\paragraph{In scope.}
(i)~Cross-user and cross-tenant disclosure through retrieval results. (ii)~Faults in the composed
components: errors in an authorizer's filter compilation; filter-semantics bugs and gaps in stores,
such as type coercion, Unicode and case handling, and nested field names; and approximate-index
behavior under selective or misaligned filters. (iii)~Authorizer outages, slow answers and malformed
answers, which must fail closed. (iv)~Identity confusion: a retrieval tool that takes a user
identifier as an argument lets a prompt-injectable model choose whose permissions apply, a confused
deputy~\cite{hardy1988confused}; the layer takes the subject only from a verified token, following
RFC 8693 and the MCP authorization specification~\cite{jones2020oauth,mcp2026model}.
(v)~Revocation: a permission removed after indexing must stop admitting chunks within the $\delta$ the
serving path states. Context added after the vector stage, such as parent documents, neighboring
chunks or graph hops, is covered only if it passes through the same re-check~\cite{thornton2026retrieval}.

\paragraph{Out of scope.}
(i)~Prompt injection carried in readable content, poisoning of readable content~\cite{zou2025poisonedrag},
and anything the model does with retrieved content, including exfiltration after correct retrieval
as in EchoLeak~\cite{nvd2025cve202532711}; we report that the model obeyed every injected notice it
retrieved (Section~\ref{sec:e2e}) to mark where our claim ends. (ii)~Wrong upstream permissions:
over-broad ACLs are enforced faithfully, as by assistants that inherit source
permissions~\cite{microsoft2026data}. (iii)~Extraction or membership inference over readable
content~\cite{zeng2024good,naseh2025riddle}. (iv)~Infrastructure: compromise of the store, its
credentials, backups or hosts. Because embeddings can be inverted~\cite{morris2023text,zhang2025universal},
our guarantee covers the query path only; the index needs storage-level protection, and the layer
never returns or logs raw vectors, query embeddings included. (v)~Several participants per
interaction, where content must be authorized for each~\cite{bhatt2025enterprise}; we model one
subject per request.

\paragraph{Side channels.}
Observation~\ref{prop:reduction} says nothing about what the timing or shape of an answer reveals.
Widening rounds, latency and the flag $c$ depend on how many unreadable chunks lie near the query, and
a filtered graph traversal may depend on the unreadable nodes it passes. Hybrid scoring that mixes
BM25 or IDF-weighted sparse scores, or normalizes or fuses scores over candidates that include
unreadable chunks, reopens B\"{u}ttcher and Clarke's collection-statistics
channel~\cite{buttcher2005security}; the dense similarity of one query and one chunk does not, because
it ignores collection statistics. We do not assume these channels away: Section~\ref{sec:side}
measures them, and we say that such a channel \emph{reveals} information about unreadable chunks,
reserving \emph{leak} for returned content.

\subsection{Metrics in full}\label{app:metrics}

Metrics use the policy oracle, which evaluates $P$ without calling the layer, and exact search over
$R_\tau(s)$ by exhaustive scoring. \emph{Selectivity} is the fraction of the corpus that passes a
filter ($|R_\tau(s)|/|D|$ for the permission filter), following ACORN~\cite{patel2024acorn}; some
systems use the complement~\cite{wang2021milvus}.

\paragraph{Leaks, and bounds per question and per subject.}
A \emph{leak} is a returned chunk outside $R_\tau(s)$ as the oracle decides it, or a canary string from
an unreadable document found in a hit, a prompt or an answer; canaries follow the Secret Sharer's
inserted secrets~\cite{carlini2019secret}, used as a deterministic enforcement check rather than a
memorization estimate. When no leak occurs, the unit of the bound matters. Hits are not independent
trials: they cluster by question and by subject, and a fault in the layer would affect a question or
a permission shape, not a random hit. We therefore bound the leak rate per question and per subject,
with the one-sided 95\% bound $1 - 0.05^{1/n} \approx 3/n$ for zero events in $n$ trials, the rule of
three~\cite{hanley1983if}. For the end-to-end runs of Section~\ref{sec:e2e} this gives about
$6.8 \times 10^{-4}$ per question for each configuration (4,420 questions) and $1.6 \times 10^{-2}$
per subject (the scenario's 189 people, who recur across runs), not the $3.8 \times 10^{-6}$ that
783,888 checked hits would suggest; for the injected faults of Section~\ref{sec:eng-revoke}, about
$2.8 \times 10^{-3}$ per fault episode (1,050 episodes). These bounds hold over the distribution of the
question generator and of the injected faults; they say nothing about an adaptive adversary, for
which only Observation~\ref{prop:reduction} and its trusted computing base speak. The
\emph{violation rate}, the share of store candidates the re-check rejects, measures how far the
pushed-down filter over-approximates, not the security of the layer.

\paragraph{Permission-recall@$k$.}
For a question whose subject can read something, permission-recall@$k$ is
$|H \cap T^\star_\epsilon| / \min(k, |R_\tau(s)|)$, with $T^\star_\epsilon$ as in
Definition~\ref{def:criteria} and $\epsilon = 10^{-4}$, so that ties do not penalize a store. Measured against exact search
over the readable set rather than the store's filtered set, it charges filter false denials and index
misses alike; recall against the store's own exact search under the same filter separates the two.

\paragraph{Under-retrieval and false completeness.}
The \emph{under-retrieval rate} is the share of retrievals that are not complete, returning fewer than
$\min(k, |R_\tau(s)|)$ hits, and the \emph{empty-result rate} the share returning none to subjects who
can read something. A count can detect under-retrieval, but mis-retrieval, where $k$ hits come back
but not the right ones, shows only in recall; in the enterprise scenario it, not empty results, caused
most of the recall loss (Sections~\ref{sec:quality} and~\ref{sec:repair}). We report the confusion
matrix of $c$ against the oracle; the \emph{false-completeness rate}, the share of results with
$c = 1$ that are not maximal, measures how far a count-honest flag stays from certifying maximality.

\paragraph{Latency.}
End-to-end latency per retrieval (filter compilation, search, re-check and widening) is reported at
p50, p95 and p99, each the median of three passes after warm-up, with thread count, batch mode and
cache state stated; widening rounds and authorizer calls are recorded for the side-channel analysis.

\paragraph{A measured correlation axis.}
Recall under a filter depends on selectivity and on where the readable set lies relative to the
query. Prior work calls the second property query correlation and controls it
explicitly~\cite{patel2024acorn,zhang2026vecbench,amanbayev2026filtered,lim2026revisiting}; our
``aligned'' and ``misaligned'' permission layouts are positive and negative instances of it, not a
new axis. We therefore place each (question, subject) pair on measured axes: global selectivity;
ACORN's query correlation, which compares the query's distance to the readable set with its expected
distance to a random set of equal size~\cite{patel2024acorn}; Global-Local
Selectivity~\cite{amanbayev2026filtered}; HCBGen's $\alpha$-hardness~\cite{lim2026revisiting}; the
local selectivity $\rho_m$, the share of readable chunks among the $m$ nearest unfiltered neighbors
($m \in \{100, 1000\}$); and the unfiltered rank of the first readable chunk. Because selectivity and
correlation alone are unstable proxies for difficulty~\cite{lim2026revisiting}, we report results
against several of these. Rates carry 95\% confidence intervals from a cluster bootstrap over
subjects, and comparisons between configurations are paired on the same questions.
